%% ============================================================
%%  Institutional Background and Identification Strategy
%%  Depositor Discipline in Modern Banking — v2 (top-journal revision)
%% ============================================================

\section{Institutional Background and Identification Strategy}
\label{sec:background}

\subsection{From Incurred Losses to CECL}
\label{sec:background:transition}

For most of the postwar period, U.S. banks recognized credit losses under the
\textit{incurred loss model} (ILM), codified in ASC 450-20 and ASC 310-10. The
ILM prohibited reserve recognition until a triggering event—typically a missed
payment or internal downgrade—made a loss ``probable and estimable.'' Anchored to
observable credit events that had already occurred, the standard deferred
recognition of latent risk embedded in seasoning loan portfolios. The result was
that two otherwise identical banks could carry materially different allowances
depending on portfolio vintage, borrower geography, and the timing of internal
credit reviews—none of which necessarily reflected differences in current
managerial risk appetite. The ILM's procyclical bias has been documented
extensively: provisioning lags default rates by one to two years over the cycle
\citep{BeattyLiao2011}, and delayed recognition reduces the informativeness of
public disclosures precisely when outside monitors most need timely signals
\citep{BushmanWilliams2012, BushmanWilliams2015}.

The 2007–2009 financial crisis prompted FASB to replace the incurred-loss trigger
with a forward-looking standard. In June 2016, FASB issued ASU 2016-13,
introducing the \textit{Current Expected Credit Loss} model (CECL), codified in
ASC 326. Under CECL, a bank must estimate and reserve for \textit{lifetime}
expected credit losses at origination, incorporating historical loss experience,
current conditions, and forward-looking forecasts. The magnitude of the required
reserve is a direct function of the credit-risk composition of the loan portfolio:
loans with longer remaining maturities, weaker borrower quality, or thinner
collateral coverage require larger CECL allowances relative to the ILM allowance
previously carried.

\paragraph{Adoption timeline.}
FASB phased adoption by entity type. Large SEC registrants adopted on January 1,
2020; smaller reporting companies, non-accelerated filers, and non-public entities
— which comprise our community-bank sample — adopted on January 1, 2023. Our
identification centers on the 2023 community-bank cohort. The earlier large-bank
adoption serves as a signal-validation sample, but its 2020Q1 timing coincides
with the onset of the COVID-19 shock, which contaminates any deposit-market tests
for that cohort (Section~\ref{sec:signal}).

\paragraph{Regulatory capital phase-in.}
Banking regulators permitted community banks to elect a three-year phase-in of
the day-one capital impact, spreading the reduction to risk-weighted capital ratios
over 2023–2025. The phase-in applied only to regulatory capital treatment; it did
not affect accounting recognition. A bank electing the phase-in still reported the
full CECL reserve and the full retained-earnings charge in its Q1 2023 Call
Report. We exploit variation in phase-in election as a heterogeneity test in
Section~\ref{sec:heterogeneity}.

\subsection{The Day-One Adjustment}
\label{sec:background:dayone}

At adoption, each bank recorded the difference between its new CECL-required
allowance and its existing ILM allowance as a cumulative-effect adjustment to
retained earnings on the opening balance sheet — the ``Day-One'' impact. For bank
$i$ adopting in quarter $t^{*}$:
\begin{equation}
  \Delta_i \;\equiv\;
  \mathrm{ALLL}^{\mathrm{CECL}}_{i,t^{*}} -
  \mathrm{ALLL}^{\mathrm{ILM}}_{i,t^{*}-1},
  \label{eq:dayone}
\end{equation}
with $\Delta_i < 0$ for the vast majority of adopters. Our treatment variable
scales this charge by beginning-of-period equity:
\begin{equation}
  \texttt{cecl\_equity}_{i}
  \;=\;
  \frac{\Delta_i}{\texttt{equity\_bop}_{i,t^{*}}}
  \times 100,
  \label{eq:treatment}
\end{equation}
expressed in percentage points, so that the measure is comparable across banks of
different sizes. We assign each bank its adoption-quarter value and hold it fixed
throughout the panel.\footnote{A small number of banks report their adjustment in
a quarter adjacent to the nominal 2023Q1 date due to fiscal-year heterogeneity
or filing lags. Our algorithm assigns the first quarter in the 2022Q3–2023Q4
window with a non-missing \texttt{cecl\_equity} observation as bank-specific
event time $t^{*}=0$.} For our main binary specification we define:
\begin{equation}
  \texttt{high\_cecl\_equity}_{i}
  \;=\;
  \mathbf{1}\!\left\{\texttt{cecl\_equity}_{i} \geq \tau\right\},
  \label{eq:binary}
\end{equation}
where $\tau$ is the top-decile cutoff of the cross-sectional distribution.
Results are robust to continuous-treatment specifications using the raw
\texttt{cecl\_equity} value.

\subsection{Identification Strategy}
\label{sec:background:id}

Our goal is to identify the causal effect of disclosed credit-risk information on
depositor behavior. The standard obstacle is that risk and deposit structure are
jointly determined: riskier banks carry fewer uninsured deposits in equilibrium,
so cross-sectional comparisons conflate information responses with pre-existing
sorting. The CECL Day-One adjustment addresses this problem because it reveals
risk-relevant information at a \textit{known, regulatory-mandated moment} and
with \textit{a magnitude predetermined by historical lending decisions}.

\paragraph{Exogenous timing.}
The 2023Q1 adoption date was set by FASB and federal banking regulators, not by
individual banks. Banks could neither delay adoption to conceal an unfavorable
charge nor accelerate it to preempt scrutiny, so the timing of disclosure is
orthogonal to bank-level conditions as of 2022.

\paragraph{Predetermined magnitude.}
The size of $\Delta_i$ in \eqref{eq:dayone} is determined by the cumulative
composition of a bank's loan portfolio accumulated well before the standard was
finalized in 2016 — loan vintages, maturities, collateral structures, and
borrower credit quality from prior years. Although banks anticipated CECL exposure
after 2016 and some gradually tightened underwriting, such responses compress
rather than eliminate cross-sectional variation in $\Delta_i$, and the residual
variation driven by pre-2016 portfolio composition is unlikely to correlate with
contemporaneous deposit-market conditions through any channel other than the
information disclosed at adoption.

\paragraph{Information, not liquidity.}
Equation~\eqref{eq:dayone} records a pure accounting charge: the Day-One
adjustment does not trigger a cash transfer, dividend cut, or borrower default.
It reduces reported book equity and regulatory capital ratios — absent phase-in
election — but leaves the cash flows and contractual terms of the underlying loans
unchanged. Any subsequent shift in deposit quantities or pricing therefore reflects
depositor reaction to the \textit{information} content of the disclosure, rather
than a concurrent deterioration in bank liquidity or cash-flow capacity. This
separability of information from cash flows is the central advantage of our
design relative to tests based on realized defaults, dividend omissions, or rating
downgrades, all of which confound new information with contemporaneous changes in
bank financial condition.

\paragraph{Empirical design.}
We implement a difference-in-differences estimator that exploits cross-sectional
variation in shock magnitude and the sharp adoption-quarter break:
\begin{equation}
  y_{it}
  \;=\;
  \alpha_i + \lambda_t
  + \beta \cdot \texttt{post}_{it} \times \texttt{high\_cecl\_equity}_{i}
  + X_{i,t-1}'\gamma
  + \varepsilon_{it},
  \label{eq:baseline}
\end{equation}
where $y_{it}$ is a deposit or funding-cost outcome; $\alpha_i$ and $\lambda_t$
are bank and calendar-quarter fixed effects; $\texttt{post}_{it}$ equals one at
or after the bank's adoption quarter; and $X_{i,t-1}$ is a vector of lagged
controls including log total assets and the equity ratio. We also estimate
continuous-treatment variants using \texttt{cecl\_equity} in place of the binary
indicator, and dynamic event-study specifications with quarter-relative-to-adoption
dummies over the window $|\texttt{qtrs\_since}| \leq 10$ (reference period
$k = -1$). The identifying assumption is that, conditional on fixed effects and
controls, shock magnitude is uncorrelated with differential pre-trends in
depositor behavior. We evaluate this assumption through parallel pre-trend tests
and complement it with a signal-validation exercise showing that equity markets
and CDS spreads priced the CECL information shock at large-bank adoption in 2020
(Section~\ref{sec:signal}).
